\subsection{叶状超曲面的法向与翘曲截面}

对余维一的叶片，前面的形状算子可直接与适配度规中的法向导数比较。取 $K_{ab}=\langle\nabla_{\partial_a}n,\partial_b\rangle$，则 $K_{ab}=-\langle B_{ab},n\rangle$；单位外法向球面的 $K$ 为正，而本章 $B$ 的法向分量为负。两种写法由这一等式转换。

\begin{figure}[htbp]
\centering
% Local TikZ geometry; styles inherited exclusively from the shared preamble.
\begin{tikzfig}[scale=.91]
\draw[gray!55] (-3.2,-.7) .. controls (-1.9,-1.25) and (1.5,-1.28) .. (3.2,-.7);
\draw[gray!65] (-3.1,1.0) .. controls (-1.4,.5) and (1.9,.55) .. (3.1,1.0);
\draw[thick] (0,-.7) ellipse (2.3 and .55);
\draw[thick] (0,1.05) ellipse (2.75 and .53);
\draw[gray!60,thin] (-2.3,-.7)--(-2.75,1.05) (2.3,-.7)--(2.75,1.05);
\draw[lecturearrow] (0,-.66)--(.15,.8) node[midway,right] {$n$};
\draw[lecturearrow] (.15,.8)--(1.25,.89) node[midway,above] {$N^a$};
\draw[lecturearrow] (-1.1,-.72)--(-1.1,.5) node[midway,left] {$N$};
\node[lecturelabel] at (2.7,1.85) {$\Sigma_{t+\dd t}$};
\node[lecturelabel] at (2.43,-1.56) {$\Sigma_t$};
\node[lecturelabel] at (-2.1,2.0) {$M\simeq I\times\Sigma$};
\end{tikzfig}

\caption{叶状分解中的相邻截面、单位法向 $n$、法向进动 $N$ 与切向位移 $N^a$。图中的椭圆只代表局部截面；一般叶不必具有椭圆形状。}\label{fig:ghs-foliation-normals}
\end{figure}

\cref{fig:ghs-foliation-normals} 说明了法向传播与沿叶的重新标记属于不同的几何操作；其中的曲面形状只是示意，真正进入 拉普拉斯算子 的是诱导度规与外曲率。
把一般几何量代入球面坐标，可以直接检验这些定义。
若 $h_{ab}(r,\theta)=a(r)^2\gamma_{ab}(\theta)$ 且 $\gamma$ 与 $r$ 无关，则 $K_{ab}=(a'/a)h_{ab}$，$K=(d-1)a'/a$。
代入 \cref{eq:ghs-laplace-evolving-slices}，得到
\[
 \Delta_g=\partial_r^2+(d-1)\frac{a'}a\partial_r+a^{-2}\Delta_\gamma.
\]
取 $a(r)=r$ 且 $\gamma$ 为标准球度规，才恢复 $\RR^d$ 中的球坐标公式。

实际的空间未必有径向中心。即便可以在一点附近构造测地球，这也仅给局部的“径向”参数；是否能够将径向量子化推广到整个空间，还需要分析作用量、边界态与演化核。这些问题将在下一组数学工具建立后处理。
